% Standalone manuscript: no external bibliography, figures, or input files.
% Compile twice with pdflatex lattice_covering_flows.tex.
\documentclass[11pt,a4paper]{article}
\usepackage[T1]{fontenc}
\usepackage{lmodern,microtype,amsmath,amssymb,amsthm,booktabs,array,xcolor}
\usepackage{fancyhdr,titlesec,fancyvrb,tikz}
\usepackage[left=24mm,right=24mm,top=24mm,bottom=24mm,headheight=15pt]{geometry}
\usepackage[colorlinks=true]{hyperref}
\definecolor{ink}{HTML}{17364A}
\definecolor{accent}{HTML}{006B68}
\definecolor{muted}{HTML}{626D78}
\hypersetup{linkcolor=accent,citecolor=accent,urlcolor=accent,
 pdftitle={Weighted fractional-part flows and sharp lattice covering bounds},
 pdfauthor={Shaitetsu}}
\titleformat{\section}{\sffamily\Large\bfseries\color{ink}}{\thesection}{.55em}{}
\titleformat{\subsection}{\sffamily\large\bfseries\color{accent}}{\thesubsection}{.55em}{}
\titleformat{\paragraph}[runin]{\bfseries\color{ink}}{}{0pt}{}
\newtheorem{theorem}{Theorem}[section]
\newtheorem{lemma}[theorem]{Lemma}
\newtheorem{proposition}[theorem]{Proposition}
\newtheorem{corollary}[theorem]{Corollary}
\theoremstyle{definition}\newtheorem{definition}[theorem]{Definition}
\theoremstyle{remark}\newtheorem{remark}[theorem]{Remark}
\DeclareMathOperator{\conv}{conv}
\DeclareMathOperator{\Vol}{Vol}
\DeclareMathOperator{\vol}{vol}
\DeclareMathOperator{\Surf}{Surf}
\DeclareMathOperator{\Span}{span}
\DeclareMathOperator{\inter}{int}
\DeclareMathOperator{\relint}{relint}
\newcommand{\R}{\mathbb R}\newcommand{\Z}{\mathbb Z}
\newcommand{\fr}[1]{\left\{#1\right\}}
\newcommand{\one}{\mathbf1}
\pagestyle{fancy}\fancyhf{}
\fancyhead[L]{\small\sffamily Fractional-part flows and lattice coverings}
\fancyhead[R]{\small\sffamily Geometry of numbers}
\fancyfoot[L]{\footnotesize\sffamily Research manuscript\;|\;18 November 2026}
\fancyfoot[R]{\small\sffamily\thepage}
\renewcommand{\headrulewidth}{.3pt}
\setlength{\parindent}{1.1em}\setlength{\parskip}{2.3pt}
\setlength{\emergencystretch}{2em}\raggedbottom

\begin{document}
\pdfbookmark[0]{Overview}{overview}
\noindent{\sffamily\small\bfseries\color{accent} GEOMETRY OF NUMBERS AND EXACT COVERING SYSTEMS}\par\vspace{3mm}
\noindent{\sffamily\bfseries\fontsize{23}{27}\selectfont\color{ink}\raggedright
Weighted fractional-part flows and sharp lattice covering bounds\par}
\medskip
\noindent{\large\sffamily Discrete surface area, covering minima, and an extremal counterexample}\par
\medskip\noindent Shaitetsu\quad$\cdot$\quad Research manuscript prepared for Keio\par
\noindent 18 November 2026\par

\paragraph{Abstract.}
We prove a uniform minimum estimate for weighted periodic sums of
fractional parts by balancing their continuous increase against their
downward jumps. A positive integral relation among the primitive vertex
directions of a simplex converts this estimate into a lattice covering.
This proves the original discrete surface-area inequality of
Codenotti--Santos--Schymura in every dimension, for arbitrary positive
real vertex lengths and rational vertex directions. It follows that
every non-hollow lattice $d$-polytope has covering radius at most $d/2$,
and that all covering minima of the standard terminal $n$-simplex are
$\mu_i=i/2$, as conjectured by Gonz\'alez Merino--Schymura.
For a simplex in the lattice generated by its vertices, we characterize
equality $\mu=d/2$ by exact partitions of the integers into residue
classes. An explicit partition yields a $12$-dimensional simplex of
normalized volume $30$ and covering radius $6$. A volume invariant
excludes it from the proposed direct-sum family of extremizers,
disproving the equality-classification clause of Conjecture A.
All arguments are unconditional. The elementary balance principle has
classical precedents; priority for the applications is not certified
by the literature search reported here.

\section{Original questions and results}
\label{sec:intro}
Let $\Lambda$ be a full-rank lattice in $\R^d$. For a convex body $K$,
its covering radius is
\[
 \mu(K,\Lambda)=\inf\{\rho\ge0:\rho K+\Lambda=\R^d\}.
\]
A lattice polytope is \emph{non-hollow} if its interior contains a
lattice point. For a $k$-dimensional body in a lattice $\Gamma$, use
normalized volume
\[
 \Vol_\Gamma(K)=\frac{k!\vol_k(K)}{\det\Gamma}.
\]
Thus a unimodular simplex has volume one; in dimension zero a point
has volume one. Quotient spaces are given the projected lattice and
the corresponding normalized volume.

The two affirmative targets are preexisting statements, with their
original constants and domains. Codenotti--Santos--Schymura
\cite[Conjecture C]{CSS} ask whether a simplex $S$ with $0\in\inter S$
and rational vertex directions satisfies
\begin{equation}\label{eq:C}
 \mu(S,\Lambda)\le
 \frac{\Surf_\Lambda(S)}{2\Vol_\Lambda(S)},
 \qquad
 \Surf_\Lambda(S)=\sum_{j=0}^{d}
       \Vol_{\pi_j\Lambda}(\pi_j S),
\end{equation}
where $\pi_j$ is projection along its $j$th vertex. A direction is
rational when its line contains a nonzero lattice vector.

For
\begin{equation}\label{eq:Tn}
 T_n=\conv\{-\one_n,e_1,\ldots,e_n\},
\end{equation}
Gonz\'alez Merino--Schymura conjecture
\begin{equation}\label{eq:B}
 \mu_i(T_n,\Z^n)=i/2\qquad(1\le i\le n).
\end{equation}
Here $\mu_i(K,\Lambda)$ is the infimum of the dilations $\rho$ for
which $\rho K+\Lambda$ meets every affine subspace of codimension $i$.
The conjecture occurs in Remark~4.9 of the 2016 preprint
\cite{GMS}, as Conjecture B of \cite{CSS}, and again as
Conjecture~4.2 of the January 2026 paper \cite{Krivokuca}.

\begin{theorem}[Discrete surface area]\label{thm:mainC}
Inequality \eqref{eq:C} holds for every $d\ge1$, every full-rank
lattice, and every simplex in its stated domain.
\end{theorem}

\begin{theorem}[Non-hollow coverings and covering minima]\label{thm:mainB}
Every non-hollow lattice $d$-polytope $P$ satisfies
$\mu(P,\Lambda)\le d/2$. Consequently \eqref{eq:B} holds for all
integers $n\ge1$ and $1\le i\le n$.
\end{theorem}

The sharp upper bound and its equality classification have different
outcomes. The full Conjecture A of \cite{CSS} also proposes that every
equality case is constructed from standard terminal simplices by
translations and direct sums, up to unimodular equivalence. Direct sum
means the convex hull of bodies in complementary coordinate spaces,
each containing the origin; the origin is allowed on their boundaries.

\begin{theorem}[Counterexample to the equality classification]
\label{thm:mainA}
Set
\[
 p=(5,5,3,3,3,3,2,1,1,1,1,1)^{\mathsf T},\qquad
 S_* =\conv\{-p,e_1,\ldots,e_{12}\}\subset\R^{12}.
\]
Then $0\in\inter S_*$, $\Vol_{\Z^{12}}(S_*)=30$, and
$\mu(S_*,\Z^{12})=6$. This simplex is not unimodularly equivalent
to any body in the proposed direct-sum and translation family.
Thus Conjecture A, including its equality clause, is false.
\end{theorem}

The key estimate and the covering construction are proved in
Sections~\ref{sec:flow}--\ref{sec:surface}.
Section~\ref{sec:polytope} proves the polytope bound, and
Section~\ref{sec:minima} proves the covering-minima formula.
Sections~\ref{sec:ecs}--\ref{sec:counterexample} give a full extremality
criterion for vertex-generated lattices and the counterexample.
Section~\ref{sec:scope} records provenance and the limits of the
priority claim. The source is self-contained, including its references
and an optional exact-arithmetic verification script.

\section{A uniform problem for weighted fractional parts}
\label{sec:flow}
Write $\{t\}=t-\lfloor t\rfloor\in[0,1)$. Fix positive integers
$r_0,\ldots,r_{m-1}$ and positive real weights
$c_0,\ldots,c_{m-1}$. Put
\[
 C=\sum_{j=0}^{m-1}r_jc_j,\qquad
 F_t(s)=\sum_{j=0}^{m-1}c_j\{t_j+r_js\}.
\]
The auxiliary question is to bound $\min_s F_t(s)$ uniformly over
\emph{every} phase vector $t\in\R^m$. Averaging over the period only
gives $\frac12\sum_jc_j$. The estimate below subtracts a definite
positive term determined by the jumps.

\begin{lemma}[Weighted fractional-part descent]\label{lem:flow}
For every $m\ge1$, every such $r,c$, and every $t\in\R^m$,
\begin{equation}\label{eq:flowbound}
 \min_{s\in[0,1)}F_t(s)
 \le B(r,c):=\frac12\sum_jc_j-
       \frac{\sum_jr_jc_j^2}{2\sum_jr_jc_j}.
\end{equation}
More precisely, let $E\subset\R/\Z$ be the distinct jump times and
let
\[
 I_e=\{j:t_j+r_je\in\Z\},\qquad J_e=\sum_{j\in I_e}c_j.
\]
Then the following identity holds with the post-jump, right-continuous
values $F_t(e)$:
\begin{equation}\label{eq:jumpidentity}
 \frac1C\sum_{e\in E}J_eF_t(e)
 =\frac12\sum_jc_j-\frac1{2C}\sum_{e\in E}J_e^2.
\end{equation}
\end{lemma}

\begin{proof}
The function is one-periodic and has constant slope $C$ between
consecutive jumps. Coordinate $j$ wraps exactly $r_j$ times in a
period, and its fractional part has period-average $1/2$. Thus
\begin{equation}\label{eq:flowaverages}
 \sum_{e\in E}J_e=C,\qquad
 \int_0^1F_t(s)\,ds=\frac12\sum_jc_j.
\end{equation}
At an event $e$, the left limit is $F_t(e)+J_e$. The change of the
square across that event is therefore
\[
 F_t(e)^2-(F_t(e)+J_e)^2=-2J_eF_t(e)-J_e^2.
\]
Choose a period interval whose endpoints are not events. Integrating
the derivative $2CF_t$ on all continuity intervals and adding the
jump changes gives
\begin{equation}\label{eq:squarebalance}
 0=2C\int_0^1F_t(s)\,ds
       -2\sum_{e\in E}J_eF_t(e)-\sum_{e\in E}J_e^2.
\end{equation}
This proves \eqref{eq:jumpidentity}. Positivity of the weights yields
\begin{equation}\label{eq:noncancellation}
 \sum_{e\in E}J_e^2
 \ge\sum_{e\in E}\sum_{j\in I_e}c_j^2
 =\sum_jr_jc_j^2.
\end{equation}
By \eqref{eq:flowaverages}, the left side of
\eqref{eq:jumpidentity} is a convex combination of the post-jump
values. At least one is at most $B(r,c)$.
The minimum is attained at an event because the slope elsewhere is
strictly positive. Coincident events were grouped throughout, so no
generic-position assumption on the phases is needed.
\end{proof}

\paragraph{The substantive estimate.}
The squared-jump term is positive and cannot cancel. Simultaneous
crossings contribute additional positive cross terms. In particular,
when every $c_j=1$, \eqref{eq:flowbound} gives the uniform gain
\begin{equation}\label{eq:unweighted}
 \min_s\sum_{j=0}^{m-1}\{t_j+r_js\}\le(m-1)/2.
\end{equation}
This gain, and its weighted version, are what change the geometric
covering estimate. An identity without this sign estimate would not
give the asserted bound.

\begin{proposition}[Equality for fixed phases]\label{prop:phaseeq}
For fixed $r,c,t$, equality holds in \eqref{eq:flowbound} if and only
if all coordinate crossings are distinct and every post-jump value
equals $B(r,c)$. Equivalently, for each event of coordinate $j$,
the preceding event gap has length $c_j/C$, and events are simple.
\end{proposition}
\begin{proof}
Equality requires equality both in the weighted-average estimate
and in \eqref{eq:noncancellation}. Every averaging weight $J_e/C$
is positive. Hence all averaged values must equal the minimum, and
each event must involve just one coordinate. Conversely these
conditions make \eqref{eq:jumpidentity} an equality at the claimed
minimum. Between two successive events the increase is $C$ times
the gap. Equal post-jump levels are therefore equivalent to the
stated gap formula, including the gap across the period boundary.
\end{proof}

\begin{figure}[ht]
\centering
\begin{tikzpicture}[x=8.4cm,y=.75cm,>=stealth]
 \draw[->,muted] (-.045,0)--(1.08,0) node[right]{$s$};
 \draw[->,muted] (0,0)--(0,5.5) node[above]{$F_t(s)$};
 \draw[dashed,muted] (0,3)--(1.03,3);
 \node[muted,anchor=west,font=\small] at (1.04,3){mean $3$};
 \draw[dotted,accent] (0,1.83333)--(1.03,1.83333);
 \draw[very thick,accent] (0,1.83333)--(.33333,3.83333);
 \draw[very thick,accent] (.33333,1.83333)--(.83333,4.83333);
 \draw[very thick,accent] (.83333,1.83333)--(1,2.83333);
 \draw[->,muted] (.33333,3.79)--(.33333,1.89);
 \draw[->,muted] (.83333,4.79)--(.83333,1.89);
 \draw[->,muted] (1,2.79)--(1,1.89);
 \foreach \x in {0,.33333,.83333,1}{\fill[accent] (\x,1.83333) circle(2pt);}
 \foreach \x/\y in {.33333/3.83333,.83333/4.83333,1/2.83333}{
   \draw[accent,fill=white] (\x,\y) circle(2pt);}
 \node[anchor=east,font=\small] at (-.01,1.83333){$11/6$};
 \node[below,font=\small] at (0,0){$0$};
 \node[below,font=\small] at (.33333,0){$1/3$};
 \node[below,font=\small] at (.83333,0){$5/6$};
 \node[below,font=\small] at (1,0){$1$};
\end{tikzpicture}
\caption{An exact equality example: $r=(1,1,1)$,
$c=(1,2,3)$ and $t=(0,-1/3,-5/6)$. The slope is $6$;
the three downward jumps have sizes $2,3,1$ in cyclic order.
All post-jump values are $11/6$, below the period average $3$.}
\label{fig:flow}
\end{figure}

The balance of derivatives and jumps used in
\eqref{eq:squarebalance} is an elementary periodic instance of
classical rate conservation; compare \cite[\S2, especially
equations (2.8)--(2.10)]{Miyazawa}. The proposed contribution here is
the phase-uniform estimate together with its precise lattice
realization and extremality consequences, rather than a claim to
have invented that general balance principle.

\section{From the flow to the original surface-area inequality}
\label{sec:surface}
Let $S=\conv(v_0,\ldots,v_d)$ satisfy the hypotheses of
Theorem~\ref{thm:mainC}. Write
\[
 v_j=\lambda_j u_j,\qquad \lambda_j>0,
\]
where $u_j$ is the primitive lattice vector pointing along $v_j$.
The lengths $\lambda_j$ may be arbitrary positive real numbers.
Let $a_j>0$ be the barycentric coordinates of the origin, so that
\begin{equation}\label{eq:barycentric}
 \sum_j a_j=1,\qquad \sum_j a_jv_j=0.
\end{equation}

\begin{lemma}[Primitive relation]\label{lem:relation}
There is a primitive positive integer vector
$r=(r_0,\ldots,r_d)$, unique with these properties, such that
$\sum_jr_ju_j=0$. If $c_j=1/\lambda_j$ and $C=\sum_jr_jc_j$,
then
\begin{equation}\label{eq:relationbary}
 a_j=\frac{r_jc_j}{C}.
\end{equation}
\end{lemma}
\begin{proof}
In a lattice basis the matrix with columns $u_j$ is integral of
rank $d$, and its kernel is a one-dimensional rational subspace.
Equation \eqref{eq:barycentric} gives the positive kernel vector
$(a_j\lambda_j)_j$. Clearing denominators in any rational generator
and choosing its positive sign gives $r$; dividing by the gcd makes
it primitive. The vector $(a_j\lambda_j)_j$ is proportional to $r$,
and $\sum_ja_j=1$ determines the constant, proving
\eqref{eq:relationbary}.
\end{proof}

\begin{proposition}[Covering from a positive relation]\label{prop:coverflow}
With this notation,
\begin{equation}\label{eq:coefficientcover}
 \mu(S,\Lambda)\le\frac12\sum_{j=0}^d
                        \frac{1-a_j}{\lambda_j}.
\end{equation}
\end{proposition}
\begin{proof}
Fix any $x\in\R^d$. Because the $u_j$ span the space, choose real
coefficients $t_j$ such that $x=\sum_jt_ju_j$. The integral relation
implies $x=\sum_j(t_j+sr_j)u_j$ for every $s\in\R$. The vector
\[
 z(s)=\sum_j\lfloor t_j+sr_j\rfloor u_j
\]
belongs to $\Lambda$, and
\begin{equation}\label{eq:rounding}
 x-z(s)=\sum_j\{t_j+sr_j\}u_j
       =\sum_jc_j\{t_j+sr_j\}v_j.
\end{equation}
The coefficients on the right are nonnegative. By
Lemma~\ref{lem:flow}, some event $s$ has coefficient sum at most
\[
 B(r,c)=\frac12\sum_jc_j-
             \frac12\sum_ja_jc_j
        =\frac12\sum_j\frac{1-a_j}{\lambda_j}.
\]
Since $0\in S$, a nonnegative vertex combination with coefficient
sum at most $B(r,c)$ belongs to $B(r,c)S$. Thus
$x\in z(s)+B(r,c)S$. The argument applies to every real $x$;
neither $x$ nor the vertex lengths are assumed rational.
\end{proof}

We next identify the coefficient expression with the original
geometric quantity. This is the normalized-volume identity in
\cite[Lemma 4.1]{CSS}; its proof is included to make the equivalence
and the factorial convention explicit.

\begin{lemma}[Chord and projection volume]\label{lem:volume}
For every $j$, let $\ell_j$ be the lattice length of the chord
$S\cap\R u_j$. Then
\begin{equation}\label{eq:volumeidentity}
 \ell_j=\frac{\lambda_j}{1-a_j},\qquad
 \frac{\Vol_{\pi_j\Lambda}(\pi_j S)}{\Vol_\Lambda(S)}
     =\frac1{\ell_j}=\frac{1-a_j}{\lambda_j}.
\end{equation}
\end{lemma}
\begin{proof}
The opposite facet $H_j$ meets the line $\R u_j$ at
\[
 p_j=\sum_{k\ne j}\frac{a_k}{1-a_j}v_k
     =-\frac{a_j}{1-a_j}v_j\in\relint H_j.
\]
Hence the chord joins $p_j$ to $v_j$ and has the stated lattice
length. Also $\pi_jH_j=\pi_jS$, since the projection of $v_j$ is
zero, which belongs to the projected facet. Extend the primitive
vector $u_j$ to a lattice basis. In these coordinates it is a
coordinate vector and the projected lattice is $\Z^{d-1}$.
The facet is an affine graph over its projection. A shear of
determinant one makes this graph horizontal, and the pyramid height
in the remaining coordinate is $\ell_j$. The pyramid formula gives
\[
 d!\vol_d(S)=\ell_j(d-1)!\vol_{d-1}(\pi_jS)
\]
in lattice coordinates. This is exactly
$\Vol_\Lambda(S)=\ell_j\Vol_{\pi_j\Lambda}(\pi_jS)$.
For $d=1$ the same identity uses the stipulated volume of a point.
\end{proof}

\begin{proof}[Proof of Theorem~\ref{thm:mainC}]
Substitute \eqref{eq:volumeidentity} into
\eqref{eq:coefficientcover} and sum over the vertices.
This gives \eqref{eq:C} with its original factor $1/2$.
\end{proof}

\paragraph{Sharpness and comparison family.}
For $S(\omega)=\conv\{-\omega_0\one_d,
\omega_1e_1,\ldots,\omega_de_d\}$, the primitive relation has
$r_j=1$. Our expression reduces to
\[
 B(\one,c)=\frac{\sum_{j<k}c_jc_k}{\sum_jc_j},
 \qquad c_j=\omega_j^{-1}.
\]
This agrees with the known exact covering radius for that family
\cite[Theorem 1.4]{CSS}. The advance in
Theorem~\ref{thm:mainC} is the estimate for every positive integral
relation among rational vertex directions, not the recomputation
of this special case.

\section{All non-hollow lattice polytopes}
\label{sec:polytope}
For a lattice simplex as in Section~\ref{sec:surface}, each
$\lambda_j$ is a positive integer. Therefore
\begin{equation}\label{eq:latticesimplex}
 \mu(S,\Lambda)\le\frac12\sum_j\frac{1-a_j}{\lambda_j}
 \le\frac12\sum_j(1-a_j)=\frac d2.
\end{equation}
An interior lattice point of a polytope need not lie in the interior
of a full-dimensional simplex formed from its vertices. The following
section-and-projection argument handles that issue; the reduction
is standard and appears in \cite[Lemma 2.1, Lemma 3.1 and
Corollary 3.2]{CSS}.

\begin{lemma}[Section and projection]\label{lem:section}
Let $K\subset\R^d$ be a convex body with $0\in\inter K$, and let
$L$ be a rational linear subspace. For the quotient projection
$\pi:\R^d\to\R^d/L$,
\[
 \mu(K,\Lambda)\le
 \mu(K\cap L,\Lambda\cap L)+\mu(\pi K,\pi\Lambda).
\]
\end{lemma}
\begin{proof}
Let $a$ and $b$ be numbers strictly larger than the projection and
section covering radii, respectively. Given $x$, choose $z\in\Lambda$
with $\pi(x-z)\in a\pi K$, and choose $y\in aK$ such that
$\pi y=\pi(x-z)$. Now $x-z-y\in L$, so choose
$w\in\Lambda\cap L$ with $x-z-y-w\in b(K\cap L)$.
Then $x-z-w\in aK+bK=(a+b)K$, by convexity. Taking $a,b$
down to their covering radii proves the inequality.
\end{proof}

\begin{proof}[Proof of the polytope assertion in Theorem~\ref{thm:mainB}]
Translation of a body preserves its covering radius, so translate
an interior lattice point to zero. Induct on $d$.
Choose an inclusion-minimal subset of the vertices whose convex
hull contains zero. Its convex hull $S$ is a simplex with
$0\in\relint S$: the coefficients of a minimal convex
representation are positive, and any affine dependence would allow
one to vary them until one became zero. Since zero is not a vertex
of $P$, the dimension $k=\dim S$ satisfies $1\le k\le d$.

The space $L=\Span S$ is rational, and $\Lambda\cap L$ is a
full-rank lattice in $L$. By \eqref{eq:latticesimplex} and inclusion,
\[
 \mu(P\cap L,\Lambda\cap L)\le
 \mu(S,\Lambda\cap L)\le k/2.
\]
If $k=d$ we are done. Otherwise $\pi P$ is a lattice polytope of
dimension $d-k$, and zero is in its interior because linear
surjections are open. The induction hypothesis gives
$\mu(\pi P,\pi\Lambda)\le(d-k)/2$.
Lemma~\ref{lem:section} adds the two bounds. Dimension one is
included in the argument, since then $k=d=1$.
\end{proof}

\section{Every covering minimum of the terminal simplex}
\label{sec:minima}
We use the standard projection description of covering minima,
introduced by Kannan--Lov\'asz \cite{KL}:
\begin{equation}\label{eq:projectionminima}
 \mu_i(K,\Lambda)=
 \max_{\substack{\pi\text{ rational projection}\\\operatorname{rank}\pi=i}}
       \mu(\pi K,\pi\Lambda).
\end{equation}
Rationality here is relative to $\Lambda$; equivalently the kernel
is spanned by lattice vectors. Formula
\eqref{eq:projectionminima} accounts for \emph{all} affine subspaces
in the definition of $\mu_i$, not just rational affine subspaces.
It is recalled explicitly in \cite[\S1.1]{Krivokuca}.

For clarity, the sharp lower bound used below is proved directly.

\begin{lemma}[An explicit deepest point of $T_d$]\label{lem:terminalsharp}
For every $d\ge1$, $\mu(T_d,\Z^d)=d/2$. The point
\[
 x=(1/(d+1),2/(d+1),\ldots,d/(d+1))
\]
is not covered by $\rho T_d+\Z^d$ for any $\rho<d/2$.
\end{lemma}
\begin{proof}
The upper bound follows from the preceding section. If
$x-z\in\rho T_d$ for some $z\in\Z^d$, there are nonnegative
$b_0,\ldots,b_d$ with sum $\rho$ such that
$x_j-z_j=b_j-b_0$. Consequently
\[
 \rho=\sum_{j=0}^db_j
 \ge\sum_{j=0}^d\{b_0+j/(d+1)\}
 =\frac d2+\{(d+1)b_0\}\ge\frac d2.
\]
For completeness, for an integer $m\ge1$ and real $y$, the multiset
$\{\{y+j/m\}:0\le j<m\}$ is
$\{(u+j)/m:0\le j<m\}$, where $u=\{my\}$.
Summing proves the fractional-part identity, including $u=0$.
\end{proof}

\begin{proof}[Proof of the covering-minima assertion in Theorem~\ref{thm:mainB}]
Every rational rank-$i$ projection of $T_n$ is a non-hollow lattice
$i$-polytope in its projected lattice. The polytope bound and
\eqref{eq:projectionminima} give $\mu_i(T_n,\Z^n)\le i/2$.
Projection onto the first $i$ coordinates sends $T_n$ onto $T_i$
and $\Z^n$ onto $\Z^i$. Extra vertices project to zero, already
inside $T_i$. Lemma~\ref{lem:terminalsharp} and
\eqref{eq:projectionminima} give the reverse bound.
\end{proof}

The result also gives $\mu_i(P,\Lambda)\le i/2$ for every
non-hollow lattice polytope $P$ and every admissible $i$, by applying
the same projection argument. The terminal simplex attains every
one of these bounds simultaneously.

\section{Extremal simplices and exact residue-class partitions}
\label{sec:ecs}
For this section, $S=\conv(v_0,\ldots,v_d)$ is a lattice simplex
with $0\in\inter S$, and the ambient lattice is assumed to be
\begin{equation}\label{eq:vertexlattice}
 \Lambda=\sum_{j=0}^d\Z v_j.
\end{equation}
Let $q=(q_0,\ldots,q_d)$ be the primitive positive integer relation
$\sum_jq_jv_j=0$, and put $Q=\sum_jq_j$.
The vertices in this section need not be primitive individually.

\begin{lemma}[Exact distance formula]\label{lem:gauge}
Let $V$ have columns $v_j$. If $x=Vt$, the distance in the gauge
of $S$ from $x$ to $\Lambda$ is
\begin{equation}\label{eq:gauge}
 \rho_S(x):=\inf_{z\in\Lambda}\inf\{\rho\ge0:x-z\in\rho S\}
   =\min_{s\in[0,1)}\sum_{j=0}^d\{t_j+sq_j\}.
\end{equation}
Moreover $\mu(S,\Lambda)=\max_x\rho_S(x)$, and this maximum is
attained on the compact torus $\R^d/\Lambda$.
\end{lemma}
\begin{proof}
Every $z\in\Lambda$ is $Vn$ for some $n\in\Z^{d+1}$.
All coefficient vectors representing $x-Vn$ have the form
$t-n+sq$, since $\ker V=\R q$. The gauge is the infimum of their
coefficient sums subject to nonnegativity. For a fixed $s$, this
requires $n_j\le\lfloor t_j+sq_j\rfloor$ for every $j$.
The largest permitted $n_j$ minimize the sum and leave exactly
the fractional parts. Interchanging the two infima proves the
equality; the periodic piecewise increasing function attains its
minimum at an event.

Because $S$ contains a neighborhood of zero, its gauge is finite,
continuous and bounded by a constant times Euclidean norm. Its
triangle inequality implies a uniform Lipschitz bound for
$\rho_S$ (apply the gauge bound in both directions). The latter
is lattice-periodic. It therefore attains its maximum on a compact
fundamental parallelepiped. The covering-radius definition gives
the asserted maximum formula.
\end{proof}

\begin{definition}
An exact covering system of $\Z$ is a finite family of residue
classes $a_j+m_j\Z$ which partition $\Z$: every integer belongs
to exactly one class. Repeated moduli are permitted.
\end{definition}

\begin{theorem}[Exact extremality criterion]\label{thm:ecs}
Under \eqref{eq:vertexlattice}, $\mu(S,\Lambda)=d/2$ if and only
if each $q_j$ divides $Q$ and there are integers $a_j$ such that
\begin{equation}\label{eq:ecs}
 \{a_j+m_j\Z:0\le j\le d\},\qquad m_j=Q/q_j,
\end{equation}
is an exact covering system.
\end{theorem}
\begin{proof}
Suppose first $\mu(S,\Lambda)=d/2$. By Lemma~\ref{lem:gauge}
choose a maximizing $x=Vt$, and write
$F(s)=\sum_j\{t_j+sq_j\}$. Its minimum is $d/2$.
Apply Proposition~\ref{prop:phaseeq} with all weights equal to one.
Every crossing is simple, every post-jump value is $d/2$, and
successive events are separated by exactly $1/Q$.
After a common rotation the event set is
$\{k/Q:0\le k<Q\}$ on the circle.
Coordinate $j$ has crossings spaced by $1/q_j$. Thus
$m_j=Q/q_j$ is an integer, and the indices of its events are one
residue class modulo $m_j$. The simple crossings partition the
$Q$ grid positions. Since each $m_j$ divides $Q$, those residue
classes partition all of $\Z$.

Conversely, assume \eqref{eq:ecs} and set $t_j=-a_j/m_j$.
The crossings of coordinate $j$ occur exactly at $s=k/Q$ with
$k\equiv a_j\pmod{m_j}$. The exact cover implies that all $Q$
events in a period are simple and equally spaced. With slope $Q$
and a downward jump of one after every gap $1/Q$, all post-jump
values are equal. Their common value is $d/2$, either by
\eqref{eq:jumpidentity} or by subtracting $1/2$ from the period
mean $(d+1)/2$. For $x=Vt$, Lemma~\ref{lem:gauge} gives
$\rho_S(x)=d/2$. The global upper bound supplies equality for the
covering radius.
\end{proof}

\begin{remark}[Why the ambient lattice matters]
If the vertices generate only a proper sublattice of the ambient
lattice, the fractional-part construction still proves an upper
bound. The lower-bound equality in \eqref{eq:gauge} need not hold,
since additional lattice translates may cover the proposed deepest
point. Theorem~\ref{thm:ecs} includes
\eqref{eq:vertexlattice} as an essential hypothesis.
\end{remark}

\section{A 12-dimensional counterexample to Conjecture A}
\label{sec:counterexample}
We give both the exact lower-bound witness and a separate invariant
which excludes the conjectured equality family.

\subsection{An exact covering certificate}
Use the following residue classes, also recorded in
Schnabel \cite[Example 3.1]{Schnabel}:
\begin{equation}\label{eq:schnabel}
\begin{split}
 &2\pmod6,\quad4\pmod6,\quad
 1\pmod{10},\quad3\pmod{10},\quad7\pmod{10},\quad9\pmod{10},\\
 &0\pmod{15},\quad
 5,6,12,18,24,25\pmod{30}.
\end{split}
\end{equation}
The residue classes themselves are prior work, not a new construction.
Their exact-cover property can be verified without an external result:
the following rows are disjoint and together list every residue modulo
$30$.
\begin{center}
\begin{tabular}{@{}ll@{}}
\toprule
Classes & Covered residues modulo $30$\\
\midrule
$2,4\pmod6$ & $2,4,8,10,14,16,20,22,26,28$\\
$1,3,7,9\pmod{10}$ & $1,3,7,9,11,13,17,19,21,23,27,29$\\
$0\pmod{15}$ & $0,15$\\
Six classes modulo $30$ & $5,6,12,18,24,25$\\
\bottomrule
\end{tabular}
\end{center}

For $S_*$ from Theorem~\ref{thm:mainA}, order the primitive relation
as
\[
 (q_0,q_1,\ldots,q_{12})=(1,5,5,3,3,3,3,2,1,1,1,1,1),
 \qquad Q=30,
\]
corresponding to $v_0=-p$ and $v_j=e_j$ for $j\ge1$.
The positive barycentric coordinates are $q_j/30$, so
$0\in\inter S_*$. The standard basis vertices generate
$\Z^{12}$. The moduli $30/q_j$ are precisely those in
\eqref{eq:schnabel}, so Theorem~\ref{thm:ecs} gives
\begin{equation}\label{eq:six}
 \mu(S_*,\Z^{12})=6.
\end{equation}

Here is a completely specified witness for the lower bound. Order
the residues as
\[
 a=(5,2,4,1,3,7,9,0,6,12,18,24,25)
\]
and put $t_j=-a_jq_j/30$. Then $x_*=\sum_jt_jv_j$ is
\begin{equation}\label{eq:deephole}
 x_* =\left(\frac12,\frac16,\frac25,\frac15,-\frac15,-\frac25,
 \frac13,-\frac1{30},-\frac7{30},-\frac{13}{30},
 -\frac{19}{30},-\frac23\right).
\end{equation}
There is one coordinate crossing at each $s=k/30$. Thus the
fractional-part function satisfies the exact all-real identity
\begin{equation}\label{eq:certifiedflow}
 \sum_{j=0}^{12}\{t_j+q_js\}=6+\{30s\}\qquad(s\in\R).
\end{equation}
Indeed both sides have slope $30$, a unit downward jump at every
$k/30$, and the same period mean $13/2$. Lemma~\ref{lem:gauge}
now shows that \emph{every} lattice translate has gauge distance
at least six from $x_*$. The certificate therefore rules out all
smaller dilations globally; it is not a finite search over lattice
points or a floating-point computation.

\subsection{An obstruction from normalized volume}
\begin{lemma}[Volume of a direct sum]\label{lem:sumvolume}
Let $K\subset V$ and $L\subset W$ be convex bodies containing zero,
where $V,W$ are complementary spaces of positive dimensions $k,l$
with lattices $\Lambda,\Gamma$. Then
\[
 \Vol_{\Lambda\oplus\Gamma}(K\oplus L)
 =\Vol_\Lambda(K)\Vol_\Gamma(L).
\]
The formula also holds when zero lies on either boundary.
\end{lemma}
\begin{proof}
Use lattice coordinates in each factor. Let $g_K$ be the extended
gauge of $K$, taking value $+\infty$ outside the cone it spans.
For $0\le t\le1$, its sublevel set is $tK$ and has volume
$t^k\vol_k(K)$. The fiber of $K\oplus L$ above $x\in K$ is
$(1-g_K(x))L$. This follows from
$g_K(x)+g_L(y)\le1$, which characterizes membership in the direct
sum, also when a gauge is extended-valued. Fubini's theorem and
the sublevel-volume formula give
\begin{align*}
 \vol_{k+l}(K\oplus L)
 &=\vol_l(L)\int_K(1-g_K(x))^l\,dx\\
 &=\vol_k(K)\vol_l(L)\,k\int_0^1t^{k-1}(1-t)^l\,dt\\
 &=\frac{k!l!}{(k+l)!}\vol_k(K)\vol_l(L).
\end{align*}
Multiplying by $(k+l)!$ in lattice coordinates proves the formula.
The calculation uses only $0\in K,L$, not interiority.
\end{proof}

\begin{proposition}[Dimension and volume of the proposed extremizers]
\label{prop:terminalvolume}
Every body formed from standard terminal simplices by a finite
sequence of translations and direct sums, with the direct-sum
lattice, has
\begin{equation}\label{eq:volumeobstruction}
 d=\sum_{j=1}^s k_j,\qquad
 \Vol=\prod_{j=1}^s(k_j+1),\qquad k_j\ge1.
\end{equation}
These identities are preserved by unimodular transformations.
No body in this family has dimension $12$ and normalized volume $30$.
\end{proposition}
\begin{proof}
For a leaf $T_k$, dimension is $k$ and normalized volume is $k+1$.
Translations preserve both quantities, while direct sums add
dimensions and, by Lemma~\ref{lem:sumvolume}, multiply normalized
volumes. Induction on the construction gives
\eqref{eq:volumeobstruction}, including intermediate translations
and arbitrary parenthesization.

The multiplicative partitions of $30$ into integers at least two
are
\[
 (30),\ (2,15),\ (3,10),\ (5,6),\ (2,3,5).
\]
The corresponding values of $\sum_j((k_j+1)-1)$ are
$29,15,11,9,7$, respectively. None is $12$.
\end{proof}

\begin{proof}[Proof of Theorem~\ref{thm:mainA}]
Interiority and covering radius were proved in
\eqref{eq:six}. Using the edges from $-p$, the normalized volume is
\[
 \Vol_{\Z^{12}}(S_*)
 =\left|\det(I_{12}+p\one_{12}^{\mathsf T})\right|
 =1+\one_{12}^{\mathsf T}p=30.
\]
The rank-one determinant identity follows, for example, by noting
that $I+p\one^{\mathsf T}$ acts as the identity on
$\ker\one^{\mathsf T}$ and as multiplication by
$1+\one^{\mathsf T}p$ on the line spanned by $p$.
Proposition~\ref{prop:terminalvolume} now excludes the proposed
family, proving the counterexample.
\end{proof}

\paragraph{Logical separation of the conjectures.}
The surface-area inequality proves the upper-bound part of
Conjecture A. The equality classification does not follow from that
inequality. Although \cite[Corollary 4.3]{CSS} is labeled as an
implication from C to A, its displayed argument establishes the
upper bound. The counterexample above concerns exactly the extra
equality clause. It leaves intact both the upper bound and the
covering-minima formula proved here.

\section{Construction, provenance, and scope}
\label{sec:scope}
\subsection{A finite witness for every input point}
For the simplex problem, choose $t$ with $x=\sum_jt_ju_j$ and
enumerate the event times
\[
 E=\bigcup_{j=0}^d
 \left\{\left\{\frac{k-t_j}{r_j}\right\}:0\le k<r_j\right\}.
\]
Evaluate $F_t(e)$ at these finitely many events and select a
minimizer. The lattice vector in \eqref{eq:rounding} then covers
the prescribed point within the bound. There are at most
$\sum_jr_j$ events. For rational data the entire construction uses
exact rational arithmetic. No polynomial bound in the input bit
length is asserted. For real data the mathematical existence proof
does not depend on an exact-comparison algorithm.

The global statements do not rest on computer testing. The short
script in Appendix~\ref{app:code} checks the explicit residue
certificate, its deep point, the determinant and volume obstruction,
and illustrative instances of the jump identity. These checks are
reproducibility aids for formulas already proved symbolically.

\subsection{What was transferred from the corpus}
The supplied Verifiable Mathematical Reasoning Corpus was inspected
at revision \texttt{5887cde} (the full identifier is in the reference)
\cite{Corpus}. Two transferable themes were useful: endpoint-aware
fractional-part identities in manuscript [030], especially the lemma
``Sawtooth certificates, including the endpoints,'' and the
fractional-part dynamics developed in manuscript [092].
The equispaced identity in Lemma~\ref{lem:terminalsharp} is an
elementary instance of the first theme. Here the one-dimensional
motion is periodic and determined by a positive lattice relation,
so no probabilistic limit theorem from the corpus is assumed.

All mathematical statements needed for the new covering construction
are proved above. In particular, the validity of other research
claims in the corpus is not a premise of any theorem in this paper.
The cross-field step is to interpret post-jump fractional coefficients
as exact representatives modulo a lattice, and then to convert
equality into a congruence partition.

\subsection{Historical status and originality}
The original surface-area question was present in the March 2019
version of \cite{CSS}; its revised preprint is dated August 2021
and its journal volume is 2022. The covering-minima question was
present in May 2016 in \cite{GMS} (whose preprint lists the second
author as Matthias Henze); the journal paper appeared in 2017.
Krivoku\'ca's January 2026 paper \cite{Krivokuca} still presents
the polytope and covering-minima statements as conjectures and
develops bounds towards them.

The literature check for this manuscript used the original papers,
their references, recent covering-minima work, and searches for the
named conjectures, fractional-part covering arguments, and links
between simplex coverings and exact covering systems. No identical
resolution or essentially identical geometric proof was identified
in the material checked as of 18 November 2026. This is a bounded
search report, not a proof of universal absence or a certification
of priority. Independent specialist review remains appropriate.

Classical ingredients are credited explicitly: derivative--jump
balance \cite{Miyazawa}, projection descriptions \cite{KL},
normalized chord identities and the section reduction \cite{CSS},
and the residue system \cite{Schnabel}. The substantive proposed
advance is their phase-uniform geometric coupling, the full
rational-direction bound, and the extremality criterion and
counterexample it produces.

\begin{center}
\small
\begin{tabular}{@{}p{.08\textwidth}p{.84\textwidth}@{}}
\toprule
Gate & Evidence and status in this manuscript\\
\midrule
G1 & The targets are Conjecture C, the original covering-minima
question, and the equality clause of Conjecture A in the cited
preexisting literature.\\[3pt]
G2 & Original bodies, constants, and quantifiers are retained.
The surface-area/coefficient equivalence is proved before use.\\[3pt]
G3 & The 2019 source and January 2026 follow-up document unresolved
status in those works. A later unlocated solution cannot be excluded
by this search alone.\\[3pt]
G4 & The proofs cover all dimensions and all real phases, with
coincident events included. The counterexample has an exact global
lower-bound identity.\\[3pt]
G5 & A positive squared-jump estimate supplies a uniform improvement;
it yields the full stated domain and a previously unassumed
extremality criterion. The general balance principle is classical.\\[3pt]
G6 & No matching solution was found in the checked sources.
Historical firstness and absence of all prior duplicates are
\emph{not certified}.\\
\bottomrule
\end{tabular}
\end{center}

Theorems here do not assert the weighted covering-minima conjecture
for arbitrary $S(\omega)$, the covering-product conjecture, or the
proposed bound for a prescribed number $k>1$ of interior lattice
points. Nor is dimension $12$ asserted to be the smallest possible
dimension of a counterexample to the equality classification.

\appendix
\section{Exact-arithmetic checks}
\label{app:code}
The following Python 3 code uses only the standard library.
All tests use rational numbers. It verifies concrete certificates;
the universal assertions are supplied by the proofs in the main text.
\begin{Verbatim}[fontsize=\footnotesize,frame=single,framesep=3mm]
from fractions import Fraction as F
from math import gcd, prod
from functools import reduce

def frac(x):
    return x - x.numerator // x.denominator

def flow_check(r, c, t):
    C = sum(ri*ci for ri, ci in zip(r, c))
    events = {}
    for ri, ci, ti in zip(r, c, t):
        for k in range(ri):
            e = frac((F(k)-ti)/ri)
            events[e] = events.get(e, F(0)) + ci
    values = {
        e: sum(ci*frac(ti+ri*e)
               for ri, ci, ti in zip(r, c, t))
        for e in events
    }
    J2 = sum(j*j for j in events.values())
    R2 = sum(ri*ci*ci for ri, ci in zip(r, c))
    assert sum(events.values()) == C
    assert sum(events[e]*values[e] for e in events) + J2/2 \
           == C*sum(c)/2
    assert J2 >= R2
    B = (sum(c)-R2/C)/2
    assert min(values.values()) <= B
    return B, values

# An unequal-weight equality example, and coincident jumps.
assert flow_check([1, 1, 1], list(map(F, [1, 2, 3])),
                  [F(0), -F(1, 3), -F(5, 6)])[0] == F(11, 6)
flow_check([1, 2, 3], [F(1), F(2), F(3)], [F(0)]*3)

# The 13 congruence classes, ordered to make q_0 = 1.
classes = [(5,30), (2,6), (4,6), (1,10), (3,10), (7,10),
           (9,10), (0,15), (6,30), (12,30), (18,30),
           (24,30), (25,30)]
q = [30//m for a, m in classes]
assert sum(q) == 30 and reduce(gcd, q) == 1 and q[0] == 1
assert all(sum((k-a) % m == 0 for a, m in classes) == 1
           for k in range(30))
t = [-F(a, m) for a, m in classes]
B, values = flow_check(q, [F(1)]*13, t)
assert B == 6 and set(values.values()) == {F(6)}
assert set(values) == {F(k, 30) for k in range(30)}

# Exact coordinates of the deepest point in the manuscript.
p = q[1:]
x = [t[i]-p[i-1]*t[0] for i in range(1,13)]
assert x == [F(1,2), F(1,6), F(2,5), F(1,5), -F(1,5),
             -F(2,5), F(1,3), -F(1,30), -F(7,30),
             -F(13,30), -F(19,30), -F(2,3)]

def determinant(a):
    a = [[F(v) for v in row] for row in a]
    ans = F(1)
    for j in range(len(a)):
        k = next(k for k in range(j, len(a)) if a[k][j])
        if k != j:
            a[j], a[k] = a[k], a[j]
            ans = -ans
        pivot = a[j][j]
        ans *= pivot
        for k in range(j+1, len(a)):
            ratio = a[k][j]/pivot
            a[k] = [u-ratio*v for u, v in zip(a[k], a[j])]
    return ans

edge_matrix = [[(1 if i == j else 0)+p[i]
                for j in range(12)] for i in range(12)]
assert determinant(edge_matrix) == 30

def factor_partitions(n, least=2):
    if n >= least:
        yield [n]
    d = least
    while d*d <= n:
        if n % d == 0:
            for tail in factor_partitions(n//d, d):
                yield [d]+tail
        d += 1

parts = list(factor_partitions(30))
assert all(prod(a) == 30 for a in parts)
assert {sum(v-1 for v in a) for a in parts} == {7,9,11,15,29}
print('All exact certificate checks passed.')
\end{Verbatim}

\begin{thebibliography}{99}
\bibitem{CSS}
G.~Codenotti, F.~Santos and M.~Schymura,
\emph{The covering radius and a discrete surface area for non-hollow
simplices}, Discrete \& Computational Geometry \textbf{67} (2022),
65--111. DOI: \href{https://doi.org/10.1007/s00454-021-00330-3}
{10.1007/s00454-021-00330-3}.
\href{https://arxiv.org/abs/1903.02866}{arXiv:1903.02866};
v1, 7 March 2019; v2, 27 August 2021.

\bibitem{GMS}
B.~Gonz\'alez Merino and M.~Schymura,
\emph{On densities of lattice arrangements intersecting every
$i$-dimensional affine subspace}, Discrete \& Computational Geometry
\textbf{58} (2017), 663--685.
\href{https://arxiv.org/abs/1605.00443}{arXiv:1605.00443};
v1, 2 May 2016, under B.~Gonz\'alez Merino and M.~Henze.
The covering-minima conjecture is Remark~4.9 in v1
(Remark~4.10 in the journal numbering cited by \cite{CSS}).

\bibitem{KL}
R.~Kannan and L.~Lov\'asz,
\emph{Covering minima and lattice-point-free convex bodies},
Annals of Mathematics (2) \textbf{128} (1988), 577--602.

\bibitem{Krivokuca}
K.~Krivoku\'ca,
\emph{Upper bounds on covering minima of convex bodies},
\href{https://arxiv.org/abs/2601.15173}{arXiv:2601.15173v1},
21 January 2026.

\bibitem{Miyazawa}
M.~Miyazawa,
\emph{Rate conservation laws: a survey},
Queueing Systems \textbf{15} (1994), 1--58.
\href{https://www.rs.tus.ac.jp/miyazawa/pdf-file/Miya1994-QS-RCL.pdf}
{Author-hosted full text}.

\bibitem{Schnabel}
O.~Schnabel,
\emph{On the reducibility of exact covering systems},
Integers \textbf{15} (2015), Paper A34.
\href{https://arxiv.org/abs/1402.3957}{arXiv:1402.3957v2};
\href{https://math.colgate.edu/~integers/p34/p34.pdf}{journal full text}.
See Example~3.1 for the residue system used here.

\bibitem{Corpus}
QEEDDD / Keio,
\emph{Verifiable Mathematical Reasoning Corpus},
\href{https://github.com/QEEDDD/Verifiable-Mathematical-Reasoning-Corpus}
{GitHub repository}, revision
\texttt{5887cdea4898f0a3aa5bf7a8e0829697bd2b0f4b},
consulted 18 November 2026.
Relevant manuscripts: [030], \emph{Arithmetic and Observation:
Classifications, Recovery, and Transcendence Barriers}, lemma labeled
\texttt{p10:gn:sawtooth}; and [092],
\emph{The discrepancy of $\{n\alpha+\beta\}$ at a fixed rational
starting point: Kesten's question and Cauchy laws with arithmetic scales}.
\end{thebibliography}
\end{document}
